% ===========================================================================
\chapter{The map as data}\label{ch:mapfile}

A digital map is not a picture. It is a database of things (roads,
rivers, lakes, buildings, places) with their positions, and the pictures
are drawn from that database whenever they are needed, at whatever zoom
and in whatever style. That is how one map can be drawn for a car's
dashboard and for a hiker's phone, and why the same data can answer
``how do I get there?''.

\section{Nodes and ways}

OpenStreetMap, the free map of the world that anyone can edit, stores its
whole database with just three kinds of element:
\begin{itemize}[nosep]
  \item a \textbf{node} is a single point, with a latitude and a longitude;
  \item a \textbf{way} is an ordered list of nodes: a road, a river, or,
    if its last node is its first, the outline of an area;
  \item a \textbf{relation} groups other elements, such as all the roads
    of a bus route.
\end{itemize}
Every element also carries \emph{tags} that say what it is, such as
\code{highway=residential} or \code{name=Mill Lane}.

Our map file is a small cousin of that design. Here is the start of it:

\begin{shell}
# Kestrel Island: an imaginary island in the North Atlantic.
# ...
map Kestrel Island

node 1 47.053148 -34.574915
node 2 47.053148 -34.573595
node 3 47.053148 -34.572275
...
road street two-way Quay Street: 1 2 3 4 5 6 7 8 9
road street one-way Market Street: 10 11 12 13 14 15 16 17 18
...
road main two-way Coast Road: 48 49 50 51 52 53 54 55 56 ...
road path two-way Lake Walk: 100 144 145 146 147 148 149 ...
...
area land Kestrel Island: 174 175 176 177 178 179 180 ...
area water Mirror Lake: 238 239 240 241 242 243 244 245 ...
river Kestrel Brook: 256 257 258 259 260 261 262 263 264 265
...
place The Harbour: 266
place Lighthouse: 271
\end{shell}

Each line is one record. A \code{node} line gives a node's id and its
position. A \code{road} line gives the road's kind (\code{main},
\code{lane}, \code{street}, \code{track} or \code{path}), whether it is
one-way, its name, and then, after the colon, its nodes in order. An
\code{area} line gives the kind of area (\code{land}, \code{water},
\code{wood}, \code{sand} or \code{town}), its name, and the nodes around its
outline. A \code{river} is a line of nodes, from its source to the sea. A
\code{place} is a name for a single node.

Keeping the positions in nodes, and having everything else refer to nodes
by id, has a lovely consequence: \emph{roads that share a node are
joined}. Node 1 is the first node of both Quay Street and Anchor Street,
so that corner is a junction, and anyone on Quay Street can turn into
Anchor Street there. Two roads that merely cross on the map without
sharing a node, such as a road on a bridge over another, are not joined.
The connections that route planning needs are all in the data, without any
geometry.

The order of a one-way road's nodes matters too: Market Street may be
driven only from node 10 towards node 18, which on the map is from west to
east.

\section{Kinds of road}

A road's kind decides how it is drawn and how fast you can travel on it. A
fixed set of named values like this is what an \emph{enumeration} is for.
In Nex, an \code{enum union} declares one: the values are
\code{Road_Kind.Main}, \code{Road_Kind.Lane} and so on, each of them a
single shared object, and \code{Road_Kind.values} is an array of them all,
in the order they are written.

\begin{newfile}{island_map.nex}
-- island_map.nex: the map itself. Nodes, roads, areas, a river and named
-- places, read from a .map file.

intern Geo
intern Geometry
intern io/Path

-- The kinds of road, from fastest to slowest. A footpath is for walkers
-- only.
enum union Road_Kind
  Main
  Lane
  Street
  Track
  Footpath
end

-- The kind written in a .map file: "main", "lane", "street", "track" or
-- "path".
function road_kind(word: String): Road_Kind do
  if word = "main" then
    result := Road_Kind.Main
  elseif word = "lane" then
    result := Road_Kind.Lane
  elseif word = "street" then
    result := Road_Kind.Street
  elseif word = "track" then
    result := Road_Kind.Track
  elseif word = "path" then
    result := Road_Kind.Footpath
  else
    raise "unknown road kind: " + word
  end
end
\end{newfile}

The file says \code{path}, but the value is called \code{Footpath}, because
\code{Path} is already the name of the library class for file paths, which
we use below to read the file.

The \code{raise} at the end matters. \code{Road_Kind} is not an optional
type, so the function must return one on every path, and there is no
sensible kind for a word we don't know. Raising an error is the honest
answer, and Nex accepts it, since a path that raises never returns.

\section{Roads, areas, rivers and places}

Each kind of record becomes a small class. They hold node \emph{ids}, not
positions, exactly as the file does.

\begin{append}{island_map.nex}
-- A road: a line through two or more nodes, in order.
class Road
create
  make(road_name: String, road_kind: Road_Kind, one_way: Boolean, ids: Array[Integer])
  require
    at_least_two_nodes: ids.length >= 2
  do
    name := road_name
    kind := road_kind
    is_one_way := one_way
    node_ids := ids
  end
feature
  once name: String
  once kind: Road_Kind
  -- A one-way road may be driven only from its first node towards its last.
  once is_one_way: Boolean
  once node_ids: Array[Integer]
invariant
  at_least_two_nodes: node_ids.length >= 2
end

-- An area: the polygon through its nodes, the last joined back to the
-- first. The kind is "land", "water", "wood", "sand" or "town".
class Area
create
  make(area_kind, area_name: String, ids: Array[Integer])
  require
    at_least_three_nodes: ids.length >= 3
  do
    kind := area_kind
    name := area_name
    node_ids := ids
  end
feature
  once kind: String
  once name: String
  once node_ids: Array[Integer]
invariant
  at_least_three_nodes: node_ids.length >= 3
end

-- A river, as a line through its nodes from source to sea.
class River
create
  make(river_name: String, ids: Array[Integer])
  require
    at_least_two_nodes: ids.length >= 2
  do
    name := river_name
    node_ids := ids
  end
feature
  once name: String
  once node_ids: Array[Integer]
end

-- A named place: a single node with a name.
class Place
create
  make(place_name: String, id: Integer) do
    name := place_name
    node_id := id
  end
feature
  once name: String
  once node_id: Integer
end
\end{append}

A road needs at least two nodes and an area at least three, so the
constructors require it. A road with one node, or a lake with two corners,
would be a mistake in the map file, and it is far better to stop when the
file is read than to crash much later while drawing.

\section{Reading the file}

The map itself keeps every node's position twice: as latitude and longitude
from the file, and as a flat point in metres, for everything else. Both are
\code{Map}s from node ids. It reads the file line by line, skipping blank
lines and comments.

\begin{append}{island_map.nex}
class Island_Map
create
  -- Read a map from a .map file (see kestrel_island.map for the format).
  load(file: String) do
    name := "Untitled"
    positions := {}
    points := {}
    roads := []
    areas := []
    rivers := []
    places := []
    let line_number: Integer := 0
    across create Path.make(file).read_text().split("\n") as raw do
      line_number := line_number + 1
      let line: String := raw.trim()
      if line /= "" and not line.starts_with("#") then
        do
          read_line(line)
        rescue
          raise file + ", line " + line_number.to_string + ": " + exception.to_string
        end
      end
    end
    if positions.is_empty then
      raise file + ": no nodes"
    end
    projection := create Projection.make(middle_of(positions.values))
    across positions as entry do
      points.put(entry.key, projection.to_point(entry.value))
    end
  end
\end{append}

Two things happen after the reading.

First, the projection. Its centre is the middle of the map: halfway
between the most northern and most southern nodes, and halfway between
the most eastern and most western. That keeps every node within a couple
of kilometres of the centre, where the projection is most accurate.

Second, every position is projected once, and the flat points stored by
node id. From here on, nothing else in the project needs to know that the
Earth is round.

If anything goes wrong reading a line, the \code{rescue} catches the
error, and raises it again with the file name and line number in front, so
that a mistake in a 300-line file can be found at once.

Now the queries, starting with positions:

\begin{continued}{island_map.nex}
feature
  name: String
  -- Every node's position on the globe, by node id...
  positions: Map[Integer, Lat_Lon]
  -- ...and on the flat map, in metres from the map's centre.
  points: Map[Integer, Point]
  projection: Projection
  roads: Array[Road]
  areas: Array[Area]
  rivers: Array[River]
  places: Array[Place]

  -- The flat position of a node.
  point(id: Integer): Point
  require
    known_node: points.contains_key(id)
  do
    result := points.get(id)
  end

  -- The flat positions of several nodes, in order.
  points_of(ids: Array[Integer]): Array[Point] do
    across ids as id do
      result.add(point(id))
    end
  ensure
    same_length: result.length = ids.length
  end

  -- The smallest and largest x and y of all the nodes.
  min_x(): Real do
    result := extreme(true, true)
  end

  max_x(): Real do
    result := extreme(true, false)
  end

  min_y(): Real do
    result := extreme(false, true)
  end

  max_y(): Real do
    result := extreme(false, false)
  end
\end{continued}

\code{points_of} adds to \code{result} straight away: a routine returning
an \code{Array} starts with a fresh empty one.

People name places loosely, so \code{find_place} accepts any part of a
name, ignoring capitals: \code{"mill"} finds Old Mill. An exact match is
looked for first, so that a name that happens to be part of another
(Long Sands inside ``Long Sands Car Park'', say) can still be chosen.

Then a question every map app answers: what is at this point? Areas can
overlap (the lake lies inside the island), and the file lists the ones
drawn on top later, so the answer is the last area in the list that
contains the point. The \code{?Place} and \code{?Area} return types say
that the answer may be \code{nil}: no such place, or out at sea.

\begin{continued}{island_map.nex}
  -- The place called `query`, ignoring case; failing that, the first
  -- place whose name contains it ("mill" finds "Old Mill").
  find_place(query: String): ?Place do
    let wanted: String := query.to_lower()
    across places as p do
      if result = nil and p.name.to_lower() = wanted then
        result := p
      end
    end
    across places as p do
      if result = nil and p.name.to_lower().contains(wanted) then
        result := p
      end
    end
  end

  -- The area at a point, or nil out at sea. Areas listed later are drawn
  -- on top of earlier ones (the lake on top of the island), so the last
  -- one that contains the point is the one you would see there.
  area_at(p: Point): ?Area do
    across areas as a do
      if polygon_contains(points_of(a.node_ids), p) then
        result := a
      end
    end
  end
\end{continued}

Last, the private routines that do the reading. A line is split at its
colon into a \emph{head}, the words before it, and the node ids after it.
Every id must already have appeared in a \code{node} line; otherwise the
file is broken, and we say so. Then the first word says what kind of
record the line is.

\begin{continued}{island_map.nex}
private feature
  read_line(line: String) do
    let colon: Integer := line.index_of(":")
    let head: Array[String] := []
    let ids: Array[Integer] := []
    if colon >= 0 then
      head := words(line.substring(0, colon))
      across words(line.substring(colon + 1, line.length)) as w do
        let id: Integer := w.to_integer
        if not positions.contains_key(id) then
          raise "node " + w + " is used before it is defined"
        end
        ids.add(id)
      end
    else
      head := words(line)
    end
    let record: String := head.get(0)
    if record = "map" then
      name := joined(head, 1)
    elseif record = "node" then
      positions.put(head.get(1).to_integer,
        create Lat_Lon.make(head.get(2).to_real, head.get(3).to_real))
    elseif record = "road" then
      roads.add(create Road.make(joined(head, 3), road_kind(head.get(1)),
        head.get(2) = "one-way", ids))
    elseif record = "area" then
      areas.add(create Area.make(head.get(1), joined(head, 2), ids))
    elseif record = "river" then
      rivers.add(create River.make(joined(head, 1), ids))
    elseif record = "place" then
      places.add(create Place.make(joined(head, 1), ids.get(0)))
    else
      raise "unknown record: " + record
    end
  end

  -- The words of s, however many spaces separate them.
  words(s: String): Array[String] do
    across s.split(" ") as w do
      if w /= "" then
        result.add(w)
      end
    end
  end

  -- The words from index `first` on, joined back with single spaces.
  joined(ws: Array[String], first: Integer): String do
    result := " ".join(ws.drop(first))
  end
\end{continued}

A name can have spaces in it, so it is everything in the head after the
fixed words: \code{joined(head, 3)} on
\code{road street one-way Market Street} is \code{Market Street}.

Finally, the middle of a list of positions, for the projection's centre,
and the extremes of the flat points, for fitting the map on a screen:

\begin{continued}{island_map.nex}
  -- The point halfway between the most northern and southern latitudes,
  -- and the most eastern and western longitudes.
  middle_of(all: Array[Lat_Lon]): Lat_Lon do
    let south: Real := 90.0
    let north: Real := -90.0
    let west: Real := 180.0
    let east: Real := -180.0
    across all as p do
      south := south.min(p.latitude)
      north := north.max(p.latitude)
      west := west.min(p.longitude)
      east := east.max(p.longitude)
    end
    result := create Lat_Lon.make((south + north) / 2.0, (west + east) / 2.0)
  end

  -- The smallest (or largest) x (or y) of all the nodes.
  extreme(along_x, smallest: Boolean): Real do
    let first: Boolean := true
    across points.values as p do
      let v: Real := when along_x then p.x else p.y end
      if first or (smallest and v < result) or (not smallest and v > result) then
        result := v
        first := false
      end
    end
  end
end
\end{continued}

\section{Exploring the island}

Now we can load the island and ask it questions. One hectare is
10,000 square metres, about the size of a football pitch and a half.

\begin{tryit}{try_island.nex}
intern Island_Map
intern Helpers

let island: Island_Map := create Island_Map.load(user_path("kestrel_island.map"))
console.print_line(island.name + ": " + island.positions.size.to_string + " nodes, "
  + island.roads.length.to_string + " roads, " + island.places.length.to_string + " places")
console.print_line("centred on " + island.projection.centre.to_string())
console.print_line((island.max_x() - island.min_x()).format(0) + " m from west to east, "
  + (island.max_y() - island.min_y()).format(0) + " m from south to north")
console.print_line("")
across island.areas as a do
  console.print_line(a.name.pad_end(" ", 16) + (polygon_area(island.points_of(a.node_ids)) / 10000.0).format(1)
    .pad_start(" ", 7) + " hectares of " + a.kind)
end
console.print_line("")
across Road_Kind.values as kind do
  let total: Real := 0.0
  across island.roads as r do
    if r.kind = kind then
      total := total + line_length(island.points_of(r.node_ids))
    end
  end
  console.print_line(kind.to_string.pad_end(" ", 10) + distance_text(total).pad_start(" ", 7))
end
console.print_line("")
across island.places as p do
  if ?island.area_at(island.point(p.node_id)) as a then
    console.print_line(p.name + when a.kind = "land" then " is on " else " is in " end + a.name)
  end
end
\end{tryit}
\tryoutput{try_island}

The island is about ten square kilometres, and the program agrees with the
map: Kestrel Peak is in Hollow Wood, and the school and the market are in
Port Kestrel.

What happens if the file has a mistake in it? This script writes a small
broken map and tries to load it:

\begin{tryit}{try_bad_map.nex}
intern Island_Map
intern Helpers

let file: String := user_path("broken.map")
create Path.make(file).write_text("map Broken\nnode 1 47.0 -34.0\nnode 2 47.0 -34.1\nroad lane two-way Short Lane: 1 2 3\n")
do
  let island: Island_Map := create Island_Map.load(file)
rescue
  console.print_line(exception.to_string.replace(user_path(""), ""))
end
\end{tryit}
\tryoutput{try_bad_map}

The error names the file, the line and the problem: Short Lane uses a node
that the file never defines.
